%% wp1 布局表示：元组/形状/步幅（论文第 2 节前半）

\wph{布局表示（Layout Representation）}

\CuTe 布局（layout）是一类功能多样的对象，能够表示范围广泛的数据排列与线程排列，在抽象物理地址、将迭代顺序与存储顺序相分离等方面极具实用价值。在本节中，我们定义形状（shape）、布局与张量（tensor）的 \CuTe 表示，并说明它们与坐标（coordinate）如何相互作用。\CuTe 布局使得泛型算法只需单一实现，便可应用于任何可折叠为该算法规范形式（canonical form）的复杂数据布局。此类算法及其应用广度的示例见第~\ref{sec:algorithms}~节。

\wpsh{元组与层次元组（Tuple 与 HTuple）}

{\tt Tuple}（元组）与 {\tt HTuple}（层次元组）概念是贯穿本工作的基础数据结构。

\begin{definition}
$\texttt{Tuple}(\mathcal{T})$ 是从集合 $\mathcal{T}$ 中选取元素构成的有限有序列表。对于 $X = (X_0,X_1,\ldots,X_{n-1}) \in \texttt{Tuple}(\mathcal{X})$，我们定义如下运算：
\begin{compactitem}
\item {\bf 阶（rank）：}rank($X$)。即元组长度 $n$。
\item {\bf 访问（access）：}$X_i$。即 \texttt{Tuple} $X$ 的第 \texttt{i} 个元素，其中 $0 \leq i < \text{rank}(X)$。
\end{compactitem}
\end{definition}

\noindent $\texttt{Tuple}(\mathcal{T})$ 是元素的扁平（flat）集合，而 $\texttt{HTuple}(\mathcal{T})$ 则是“由 $\mathcal{T}$ 组成的层次元组”。

\begin{definition}
$\texttt{HTuple}(\mathcal{T})$ 或者是集合 $\mathcal{T}$ 中的一个元素，或者是一个 $\texttt{Tuple}(\texttt{HTuple}(\mathcal{T}))$。对于 $X \in \texttt{HTuple}(\mathcal{X})$，我们定义如下运算：
\begin{compactitem}
\item {\bf 阶（rank）：}rank($X$)。若 $X \in \texttt{Tuple}$，则为元组长度，否则为 1。
\item {\bf 访问（access）：}$X_i$。即 \texttt{HTuple} $X$ 的第 \texttt{i} 个元素，其中 $0 \leq i < \text{rank}(X)$。
\item {\bf 深度（depth）：}depth($X$)。若 $X \in \texttt{Tuple}$，则为 $1 + \text{max}(\text{depth}(X_0), \text{depth}(X_1), \ldots)$，否则为 0。
\end{compactitem}
\end{definition}

例如，
\begin{align*}
31 \quad\quad (16,32) \quad\quad (3,-8,7) \quad\quad (2,(4,1),-1) \quad\quad ((4,6),(3,(2,2),8))
\end{align*}
均为 $\texttt{HTuple}(\Z)$ 的实例。

%For convenience, we will also often use the following abbreviation for {\tt HTuple($T$)} operators:
%\begin{align*}
%\texttt{op<$i_0$,$i_1$,...,$i_n$>(HTuple($T$))} \Longleftrightarrow %\texttt{op(get<$i_n$>(...get<$i_1$>(get<$i_0$>(HTuple($T$))...))}
%\end{align*}

在对 \texttt{HTuple} 进行推理时，定义{\em 同余}（congruence）与{\em 弱同余}（weak congruence）的概念会很有用。
%We say that congruent \texttt{HTuple}s have the same hierarchical {\em profile}, but the value and type of each leaf element may differ.
\begin{definition}
\emph{同余} $\sim$ 是 \texttt{HTuple} 上的一个等价关系。对于 $P \in \texttt{HTuple}(\mathcal{P})$ 与 $S \in \texttt{HTuple}(\mathcal{S})$，
\begin{align*}
P \sim S \quad \text{ 当且仅当 } \quad
\begin{cases}
P \in \mathcal{P} \text{ 且 } S \in \mathcal{S}, \text{ 或} \\
P,S \in \texttt{Tuple} \text{ 且 } \text{rank}(P) = \text{rank}(S) \text{ 且 } \forall_i \ P_i \sim S_i
\end{cases}
\end{align*}
此时我们称 $P$ 与 $S$ {\em 同余}，并具有相同的{\em 轮廓}（profile）。
\end{definition}
例如，
\begin{align*}
(4,8) \sim (5,7) \quad \text{且} \quad (4,(2,4)) \sim (7,(3,2)) \quad \text{且} \quad (\v{v}, ((\v{p}, 3))) \sim (0, ((0, 0)))
\end{align*}
但 $(4,8)$ 与 $(4,(2,4))$ 不同余，$(4,(2,4))$ 与 $(0, ((0, 0)))$ 也不同余。

类似地，弱同余检验的是一个 {\tt HTuple} 的轮廓是否至少与另一个一样{\em 精细}。
\begin{definition}
\emph{弱同余} $\lesssim$ 是 \texttt{HTuple} 上的一个偏序。对于 $P \in \texttt{HTuple}(\mathcal{P})$ 与 $S \in \texttt{HTuple}(\mathcal{S})$，
\begin{align*}
P \lesssim S \quad \text{ 当且仅当 } \quad
\begin{cases}
P \in \mathcal{P}, \text{ 或} \\
P,S \in \texttt{Tuple} \text{ 且 } \text{rank}(P) = \text{rank}(S) \text{ 且 } \forall_i \ P_i \lesssim S_i
\end{cases}
\end{align*}
此时我们称 $P$ 与 $S$ {\em 弱同余}，$P$ {\em 粗化}（coarsens）$S$ 的轮廓，且 $S$ {\em 细化}（refines）$P$ 的轮廓。
\end{definition}
例如，
\begin{align*}
30 \lesssim (a,b) \lesssim (\v{v},(0,\alpha)) \quad \text{且} \quad 30 \lesssim (a,b,c) \lesssim ((0,0),0,0)
\end{align*}
但 $(a,b)$ 与 $(a,b,c)$ 不弱同余，$(\v{v},(0,\alpha))$ 与 $((0,0),0)$ 也不弱同余。

\wpsh{形状（Shape）}
\label{sec:shape}

多维数组通常由其\emph{形状}（shape）来刻画：一列描述各个模（mode）规模的正整数。$M{\times}N$ 矩阵的二维形状表示为 $(M,N)$，并自然地由坐标 $(m,n)$ 索引，其中 $0 \leq m < M$ 且 $0 \leq n < N$。对此的一个自然推广，是将形状表示为由正整数构成的层次元组。

\begin{definition}
\emph{形状}是一个 $\texttt{HTuple}(\Z^+)$，其中 $\Z^+ = \{1,2,3,\ldots\}$ 为正整数集合。形状 $S$ 的阶（rank）是该 \texttt{HTuple} 的阶。形状 $S$ 的大小（size）是其全部元素的乘积，记作 $\abs{S} = \prod_k \abs{S_k}$。
\end{definition}

层次化形状之所以特别有用，在于它们可以由多种坐标系索引。考虑不超过 $N$ 的整数集合
\begin{align*}
\Z_N = \{ 0, 1, 2, \ldots, N-1 \}.
\end{align*}
\CuTe 注意到：只要存在一个双射（bijection）
\begin{align*}
S : \Z_{MN} \longleftrightarrow \Z_M \times \Z_N
\end{align*}
在一维整型坐标（integral coordinate）$i \in \Z_{MN}$ 与二维自然坐标（natural coordinate）$(m,n) \in \Z_M \times \Z_N$ 之间建立映射，二维形状 $(M,N)$ 也可以被解释为描述 $MN$ 个一维元素，并由整型坐标 $i$（$0 \leq i < MN$）索引。

类似地，二维形状 $(M,NP)$ 可以被解释为层次形状 $(M,(N,P))$，由自然坐标 $(m,(n,p))$ 索引，其中 $0 \leq m < M$、$0 \leq n < N$、$0 \leq p < P$。相应的双射为
\begin{align*}
S : \Z_M \times \Z_{NP} \longleftrightarrow \Z_M \times (\Z_N \times \Z_P)
\end{align*}
在二维坐标 $(m,q) \in \Z_M \times \Z_{NP}$ 与自然坐标 $(m,(n,p)) \in \Z_M \times (\Z_N \times \Z_P)$ 之间建立映射。

层次形状与坐标的一个直接推论是：张量算法可以按照对它们而言最自然的形状来编写（第~\ref{sec:algorithms}~节）：{\tt COPY} 中的向量用一维形状，{\tt GEMM} 中的矩阵用二维形状，{\tt batched-GEMM} 中的张量用三维形状，等等，同时仍能接受被折叠（fold）成与算法规约弱同余的层次化形状张量（第~\ref{sec:compatibility}~节）。数据张量的形状通常表示为一列扁平的整数，它可以被任意折叠成泛型张量算法所接受的形状。另外，由于张量的每个模都关联着一个用于索引数据的步幅（stride）（第~\ref{sec:stride}~节），这种模的折叠使得我们能够表示远比 {\tt COPY} 中的简单连续数组或 {\tt BLAS GEMM} 中的行主序、列主序矩阵复杂得多的数据布局（第~\ref{sec:layout}~节）。

下一小节定义这种相容性（compatibility）的概念，以及形状内各坐标集合之间的这些关系。

\wpssh{坐标集合与相容性（Coordinate Sets and Compatibility）}
\label{sec:compatibility}

层次形状支持由多种坐标系索引。这里，我们为特定形状定义坐标集合（coordinate set），并定义形状之间的相容性概念，以便在形状之间共享坐标集合。

\begin{definition}
\emph{坐标集合}是一个由非负整数构成的集合 $\Z_N = \{0, 1, 2, \ldots, N-1\}$，或坐标集合的笛卡尔积 $\Z_N \times \Z_M = \Z_{(N,M)}$。
\end{definition}

下面是几个坐标集合：
\begin{align*}
\Z_6 &= \{ 0, 1, 2, 3, 4, 5 \} \\
\Z_3 \times \Z_4 = \Z_{(3,4)} &= \{(0,0), (1,0), (2,0), (0,1), (1,1), (2,1), (0,2), (1,2), (2,2), (0,3), (1,3), (2,3) \} \\
(\Z_2 \times \Z_1) \times \Z_3 = \Z_{((2,1),3)} &= \{((0,0),0), ((1,0),0), ((0,0),1), ((1,0),1), ((0,0),2), ((1,0),2) \}
\end{align*}
坐标集合 $\Z_S$ 恰好是形状 $S$ 的自然坐标集合。形状 $S$ 的其它坐标集合，是任何与 $S$ {\em 相容}（compatible）且{\em 粗化}（coarsens）$S$ 的形状的坐标集合。

\begin{definition}
\emph{相容性} $\preceq$ 是形状集合上的一个偏序。对于形状 $P$ 与 $S$，
\begin{align*}
P \preceq S \quad \text{ 当且仅当 } \quad
\begin{cases}
P \in \Z^+ \text{ 且 } P = \abs{S}, \text{ 或} \\
P,S \in \texttt{Tuple} \text{ 且 } \text{rank}(P) = \text{rank}(S) \text{ 且 } \forall_i \ P_i \preceq S_i
\end{cases}
\end{align*}
此时我们称 $P$ 与 $S$ {\em 相容}，$P$ {\em 粗化} $S$，且 $S$ {\em 细化} $P$。
\end{definition}

相容性要求两个形状大小相同，因此 \texttt{HTuple} 中的整数值是有影响的。例如，
\begin{align*}
30 \preceq (2,15) \preceq (2,(3,5)) \quad \text{且} \quad
30 \preceq (6,5) \preceq ((3,2),5)
\end{align*}
但 $(2,(3,5))$ 与 $((3,2),5)$ 尽管大小相同，却并不相容。不过，它们确实共享一个公共的相容形状 $30$。

有了坐标集合与形状相容性的定义，我们便可以为任意给定形状定义全部相容坐标构成的集合。

\begin{definition}
形状 $S$ 定义了一个{\em 相容坐标集合}的集合 $\Z(S)$，即所有粗化 $S$ 的形状的坐标集合。
\begin{align}
\Z(S) = \{\Z_{S'} \, \mid \, S' \preceq S\}.
\end{align}
\end{definition}

每个形状都有一个整型坐标集合，
\begin{align*}
\{ 0, 1, 2, \ldots, \abs{S}-1 \} = \Z_{\abs{S}} \in \Z(S),
\end{align*}
且每个 $r$ 阶形状都有一个 $r$ 阶坐标集合，
\begin{align*}
\{ (a_0, \ldots, a_{r-1}) \ \mid \ a_i \in \Z_{\abs{S_i}} \} = \Z_{(\abs{S_0}, \abs{S_1}, \ldots, \abs{S_{r-1}})} \in \Z(S).
\end{align*}

注意，若形状 $P$ 粗化形状 $S$，则 $\Z(P) \subseteq \Z(S)$。这意味着形状 $P$ 内的任何坐标也是形状 $S$ 内的坐标。

% {\bf Reflexivity:} For all shapes $A$, $A \preceq A$.

% {\bf Anti-symmetry:} For all shapes $A$ and $B$, if $A \preceq B$ and $B \preceq A$, then $A = B$.

% {\bf Transitivity:} For all shapes $A$ and $B$ and $C$, if $A \preceq B$ and $B \preceq C$, then $A \preceq C$.

% {\bf Minimal Elements:} For all integers $n$ and all shapes $S$ with $n = \abs{S}$, we have $n \preceq S$.

% {\bf Maximal Elements:} If all elements of a shape $P$ are prime, then there is no shape $A \neq P$ with $P \preceq A$.

\wpssh{坐标（Coordinates）}

在本小节中，我们定义若干类坐标，在形状的各相容坐标集合之间定义双射，并给出这些坐标映射的示例。

\begin{definition}
形状 $S$ 的\emph{界内坐标}（in-bounds coordinate），或简称{\em 坐标}，是其某个坐标集合中的一个元素，即 $c \in \Z_{S'} \in \Z(S)$。注意，坐标总是一个 $\texttt{HTuple}(\N)$。当意图明确时，我们将简写为 $c \in \Z(S)$。
\end{definition}

\begin{definition}
形状 $S$ 的\emph{整型坐标}（integral coordinate）是坐标 $\bar{c} \in \Z_{\abs{S}} \in \Z(S)$。注意，整型坐标总是一个整数，即 $\bar{c} \in \N$。
\end{definition}

\begin{definition}
形状 $S$ 的\emph{自然坐标}（natural coordinate）是坐标 $\tilde{c} \in \Z_S \in \Z(S)$。注意，自然坐标总是一个与形状同余的 $\texttt{HTuple}(\N)$，即 $\tilde{c} \sim S$。
\end{definition}

为了在界内坐标之间进行转换，我们在形状 $S$ 的各坐标集合上构造一个枚举，以定义{\em 坐标列表}（coordinate list）。在本工作中，我们选择坐标的余字典序（colexicographical ordering）$<$，其定义如下：
\begin{align*}
(a_0, \ldots, a_n) < (b_0, \ldots, b_n)
\quad \text{当且仅当} \quad
\begin{cases}
    a_n < b_n, \text{ 或} \\
    a_n = b_n \ \ \text{且} \ \ (a_0,...,a_{n-1}) < (b_0,...,b_{n-1})
\end{cases}
\end{align*}
并在需要时递归应用。余字典序枚举在坐标列表上定义了一个双射。函数
% \begin{align}
% \texttt{idx2crd} \colon \ \Z_{\abs{S}} &\to \Z_{(\abs{S_0}, \abs{S_1}, \ldots, \abs{S_{r-1}})}, \nonumber \\
% i &\mapsto \Bigl(i \bmod \abs{S_0}, \floor{\frac{i}{\abs{S_0}}} \bmod \abs{S_1}, \ldots, \floor{\frac{i}{\abs{S_0} \cdot \ldots \cdot \abs{S_{r-3}}}} \bmod \abs{S_{r-2}}, \floor{\frac{i}{\abs{S_0} \cdot \ldots \cdot \abs{S_{r-2}}}} \Bigr)
% \label{eqn:idx2crd}
% \end{align}
\begin{align}
\texttt{idx2crd} \colon \ \Z_{\abs{S}} &\to \Z_{(\abs{S_0}, \abs{S_1}, \ldots, \abs{S_{r-1}})}, \nonumber \\
i &\mapsto \Bigl(i \bmod \abs{S_0}, \floor{\frac{i}{\abs{S_0}}} \bmod \abs{S_1}, \ldots, \floor{\frac{i}{\prod_{k=0}^{r-3} \abs{S_k}}} \bmod \abs{S_{r-2}}, \floor{\frac{i}{\prod_{k=0}^{r-2} \abs{S_k}}} \Bigr)
\label{eqn:idx2crd}
\end{align}
将 $\Z_{\abs{S}}$ 的第 $i$ 个坐标（形状 $S$ 的第 $i$ 个整型坐标）映射为 $\Z_{(\abs{S_0}, \abs{S_1}, \ldots, \abs{S_{r-1}})}$ 的第 $i$ 个坐标（形状 $(\abs{S_0}, \abs{S_1}, \ldots, \abs{S_{r-1}})$ 的第 $i$ 个自然坐标）。其它双射，例如反向和/或镜像的字典序或余字典序，同样可以使用。

\texttt{idx2crd} 的逆由下式给出
\begin{align}
\texttt{crd2idx} \colon \ \Z_{(\abs{S_0}, \abs{S_1}, \ldots, \abs{S_{r-1}})} &\to \Z_{\abs{S}}, \nonumber \\
(c_0, c_1, \ldots, c_{r-1}) &\mapsto c_0 + c_1 \cdot \abs{S_0} + \ldots + c_{r-1} \cdot \prod_{k=0}^{r-2} \abs{S_k}
\label{eqn:crd2idx}
\end{align}
它将 $\Z_{(\abs{S_0}, \abs{S_1}, \ldots, \abs{S_{r-1}})}$ 的第 $i$ 个坐标映射为 $\Z_{\abs{S}}$ 的第 $i$ 个坐标。

若两个形状相容，即 $P \preceq S$，则 $\Z_P$ 中的坐标可以通过反复应用 \texttt{idx2crd} 映射到 $\Z_S$，而 $\Z_S$ 中的坐标可以通过反复应用 \texttt{crd2idx} 映射到 $\Z_P$。按照坐标的余字典序，图~\ref{fig:coordlists} 列出了每个形状从整型坐标到自然坐标的映射。

\begin{figure}[ht]
\centering
\begin{subfigure}[t]{0.2\textwidth}
    \centering
    \begin{tabular}{|c|c|}
    $\Z_4$ \\
    \hline
    0 \\
    1 \\
    2 \\
    3 \\
    \hline
    \end{tabular}
    \caption*{$S = 4$ \\ $\Z(S) = \{\Z_4\}$}
\end{subfigure}
\begin{subfigure}[t]{0.3\textwidth}
    \centering
    \begin{tabular}{|c|c|}
    $\Z_6$   & $\Z_{(2,3)}$ \\
    \hline
    0        & $(0,0)$ \\
    1        & $(1,0)$ \\
    2        & $(0,1)$ \\
    3        & $(1,1)$ \\
    4        & $(0,2)$ \\
    5        & $(1,2)$ \\
    \hline
    \end{tabular}
    \caption*{$S = (2,3)$ \\ $\Z(S) = \{\Z_6, \Z_{(2,3)}\}$}
\end{subfigure}
\begin{subfigure}[t]{0.4\textwidth}
    \centering
    \begin{tabular}{|c|c|c|}
    $\Z_{12}$ & $\Z_{(6,2)}$ & $\Z_{((2,3),2)}$ \\
    \hline
    0         & $(0,0)$         & $((0,0),0)$ \\
    1         & $(1,0)$         & $((1,0),0)$ \\
    2         & $(2,0)$         & $((0,1),0)$ \\
    3         & $(3,0)$         & $((1,1),0)$ \\
    4         & $(4,0)$         & $((0,2),0)$ \\
    5         & $(5,0)$         & $((1,2),0)$ \\
    6         & $(0,1)$         & $((0,0),1)$ \\
    7         & $(1,1)$         & $((1,0),1)$ \\
    8         & $(2,1)$         & $((0,1),1)$ \\
    9         & $(3,1)$         & $((1,1),1)$ \\
    10        & $(4,1)$         & $((0,2),1)$ \\
    11        & $(5,1)$         & $((1,2),1)$ \\
    \hline
    \end{tabular}
    \caption*{$S = ((2,3),2)$ \\ $\Z(S) = \{\Z_{12}, \Z_{(6,2)}, \Z_{((2,3),2)}\}$}
\end{subfigure}
\caption{各形状 $S$ 的坐标列表集合示例。}
\label{fig:coordlists}
\end{figure}

\paragraph{越界坐标（Out-of-bounds Coordinates）}除已定义的坐标集合之外，定义一些可能不在形状坐标集合之内、具有特定轮廓的坐标也很有用。

\begin{definition}
形状 $S$ 的\emph{可容许坐标}（admissible coordinate）是任意与该形状弱同余的坐标 $c \in \texttt{HTuple}(\Z)$，即 $c \lesssim S$。
\end{definition}

\begin{definition}
形状 $S$ 的\emph{越界坐标}（out-of-bounds coordinate）是任意不在界内的可容许坐标 $c \in {\tt HTuple}(\Z)$，即 $c \notin \Z(S)$。
\end{definition}

\begin{definition}
形状 $S$ 的\emph{同余坐标}（congruent coordinate）是任意与该形状同余的坐标 $c \in \texttt{HTuple}(\Z)$，即 $c \sim S$。将其记作 $\Z^S = \{c \in \texttt{HTuple}(\Z) \ \mid \ c \sim S\}$。
\end{definition}

也就是说，$\Z_S$ 是由形状 $S$ 界定的有限坐标集合，而 $\Z^S$ 是与 $S$ 同余的全部坐标构成的无限集合。由于对 $\Z^S$ 而言，形状 $S$ 内部的具体取值并不重要，我们有时会用 $(\ast,\ast)$ 作为占位轮廓（placeholder profile）。例如，
\begin{align*}
\Z^{(\ast,\ast)} = \{ (a,b) \ \mid \ a,b \in \Z \} \quad \text{且} \quad \Z^{(\ast,(\ast,\ast))} = \{ (a,(b,c)) \ \mid \ a,b,c \in \Z \}
\end{align*}

注意，{\tt idx2crd} 对所有整数（等价地说，对 $\Z^{\abs{S}}$ 中的坐标）都有良定义，而不仅限于 $\Z_{\abs{S}}$ 中的整数。当它在某个整数 $i \geq \abs{S}$ 上求值时，它总会返回一个相对于形状 $(\abs{S_0},\abs{S_1},\ldots,\abs{S_{r-1}})$ 越界的坐标 $(c_0,c_1,\ldots,c_{r-1})$。相比之下，{\tt crd2idx} 无法保证对越界坐标的输入给出越界的结果。因此，{\tt crd2idx} 与 {\tt idx2crd} 只有在界内坐标上求值时才互为逆。

\wpsh{步幅（Stride）}
\label{sec:stride}

上一节描述了形状、形状的层次结构以及这些形状的坐标。为了构造数据、线程或其它对象的\textbf{布局}，我们定义一个从形状内坐标到偏移（offset）的映射。

\begin{definition}
形状 $S$ 的\emph{步幅}（stride）$D$ 是一个与该形状同余的 $\texttt{HTuple}(\mathcal{D})$，即 $S \sim D$。该步幅定义了一个从自然坐标 $\tilde{c} \in \Z_S$ 到陪域（codomain）$\D$ 的映射，由下式给出：
\begin{align}
\texttt{inner\_product} \colon \ &\Z \cdot \D \to \D, \nonumber \\
&c \cdot d \mapsto cd \nonumber\\
\texttt{inner\_product} \colon \ &\texttt{HTuple($\Z$)} \cdot \texttt{HTuple($\D$)} \to \D, \nonumber \\
&c \cdot d \mapsto \sum_i \texttt{inner\_product}(c_i, d_i)
\label{eqn:innerprod}
\end{align}
\end{definition}

在大多数情况下，步幅也是 \texttt{HTuple($\Z$)}，即 $\D = \Z$。\texttt{inner\_product} 所产生的整数通常被解释为数据数组内的偏移。然而，步幅元素的概念可以推广为整数半模（integer-semimodule）中的任意元素，布局能表示的函数范围也宽得多。

\wpssh{整数半模（Integer-Semimodules）}

\begin{definition}
\emph{整数半模}（integer-semimodule）是一个配备了结合加法 $M + M \to M$ 与标量乘法 $\Z \cdot M \to M$ 的集合 $M$。对于 $a,b \in \Z$ 与 $m,n,p \in M$，加法与标量乘法满足
\begin{compactitem}
\item 乘法单位元：$1 \cdot m = m$，
\item 加法结合律：$m + (n + p) = (m + n) + p$，
\item 乘法结合律：$a \cdot (b \cdot m) = (ab) \cdot m$。
%\item Distributivity over group addition: $a \cdot (m+n) = a \cdot m + a \cdot n$,
%\item Distributivity over scalar addition: $(a+b) \cdot m = a\cdot m + b \cdot m$.
\end{compactitem}
不要求存在加法单位元与加法逆元，因此 $(M,+)$ 是一个半群（semigroup）。我们将整数半模记作 $(M,+,\cdot)$，用以表示集合 $M$、加法运算 $+$ 与标量乘法 $\cdot$。
\end{definition}

整数集 $\Z$ 是一个整数半模。有理数集 $\Q$ 是一个整数半模。由模 2 算术运算构成的域 $\mathbb{F}_2 = (\{0,1\},\text{XOR},\text{AND})$ 是一个整数半模。在逐元素加法与标量乘法之下，任意整数半模的笛卡尔积或 {\tt HTuple} 都是整数半模。

一个格外有用的整数半模是 $(\Z^S,+,\cdot)$，其中 $\Z^S$ 是与 $S$ 同余的全部 {\tt HTuple}($\Z$) 构成的集合。例如，2 阶算术元组的基元素构成一个整数半模：
\begin{align*}
e_0 = (1,0), \quad e_1 = (0,1), \quad \Z^{(\ast,\ast)} &= \{a \cdot e_0 + b \cdot e_1 \ \mid \ a, b \in \Z\}
\end{align*}
其整数缩放与加法按元素定义：
{\small
\begin{align*}
a \cdot e_0 + b \cdot e_1 = (a, b)
\end{align*}}

因此，$e_0$、$e_1$ 以及任意线性组合都可以用作布局中的步幅。通过从 $\Z^S$ 中选取步幅元素，布局可以经由 \texttt{inner\allowbreak\_product} 运算生成形状 $S$ 的自然坐标。
